\section{Rangkuman}
\label{sec:rangkuman-bab13}

Pada bab ini telah dibahas:

\begin{itemize}
    \item Label adalah nama simbolik untuk lokasi di dalam segmen
          kode. Label didefinisikan dengan titik dua dan ditulis
          tanpa titik dua saat menjadi tujuan lompatan.

    \item \asm{JMP} memindahkan kendali tanpa syarat, sedangkan
          \asm{Jcc} memindahkan kendali hanya bila syarat
          \eng{flag} terpenuhi.

    \item \asm{CMP} mengurangkan operand kedua dari operand pertama
          tanpa menyimpan hasil. Yang berubah hanya \eng{flag}.

    \item Lompatan bersyarat untuk bilangan bertanda memakai
          \asm{JG}, \asm{JGE}, \asm{JL}, dan \asm{JLE}. Untuk
          bilangan tidak bertanda dipakai \asm{JA}, \asm{JAE},
          \asm{JB}, dan \asm{JBE}.

    \item Pola \code{if} diterjemahkan menjadi \asm{CMP} diikuti
          satu lompatan bersyarat. Pola \code{if--else} memerlukan
          tambahan \asm{JMP} untuk melewati blok \code{else}.

    \item \asm{LOOP} mengurangi \reg{CX} dengan satu, lalu melompat
          bila \reg{CX} belum nol. Badan perulangan selalu dijalankan
          paling sedikit satu kali.

    \item Perulangan \code{while} memeriksa syarat sebelum badan,
          sedangkan \code{for} dengan cacah tetap dapat memakai
          \asm{LOOP}.

    \item Perulangan bersarang memerlukan penyimpanan pencacah luar,
          biasanya melalui \eng{stack} dengan \asm{PUSH} dan
          \asm{POP}.
\end{itemize}

\begin{checklistbox}
\vspace{1.1em}
Tandai setiap butir yang sudah Anda kuasai.

\begin{itemize}
    \item[\checkbox] Saya dapat menulis label dan menggunakannya
          sebagai tujuan lompatan.
    \item[\checkbox] Saya dapat menggunakan \asm{JMP}.
    \item[\checkbox] Saya dapat menggunakan \asm{CMP} untuk
          menyiapkan \eng{flag}.
    \item[\checkbox] Saya dapat memilih \asm{Jcc} untuk bilangan
          bertanda dan tidak bertanda.
    \item[\checkbox] Saya dapat menyusun pola \code{if} dan
          \code{if--else}.
    \item[\checkbox] Saya dapat menyusun perulangan dengan
          \asm{LOOP}.
    \item[\checkbox] Saya dapat menyusun perulangan bersarang untuk
          array.
\end{itemize}

\vspace{0.4em}
Bila masih ada butir yang belum bertanda, ulangi bagian yang
bersesuaian sebelum melanjutkan ke bab berikutnya.
\end{checklistbox}

\begin{infobox}
\vspace{1.1em}
\textbf{Rujukan bab.} Materi pada bab ini diringkas dari
\cite{gunawan2019}, Bagian 11, hlm.~101--114, dan
\cite{sto2001}, BAB XIII, hlm.~68--72.
\end{infobox}
